\scanpage{134}

et on choisit un \inserted{$L$-}plongement de $\tilde L_J$ dans $\bar L$.

C'est par cela que (37) est définie --- c'est le \inserted{début de la suite d'homotopie du} fibré, \unsure{moralement}, d'un produit de disques sur $D_J$ (savoir $U_J$ lui-$\unsure{=}$) --- j'ai eu un peu peur d'y mettre un zéro à g, à cause du $\deleted{\unsure{\mathcal E}}$
\[
\pi_2(\bar D_J)\longrightarrow
\underset{\substack{\Vert\\ T^J}}{\pi_1(\text{fibre})}
\qquad (?)!
\]
Mais considérons l'hom
\[
\mathcal E_{U_J}^{\tilde U_J}\longrightarrow\mathcal E_U^{\bar L}
\]
induit par ``l'injection'' $U_J\longrightarrow U$ ; \deleted{elle}\inserted{\deleted{devrait}} heuristiquement on \deleted{devrait avoir}\inserted{a} un diagramme \deleted{induit par} commutatif
\diagram{
\begin{tikzcd}[row sep=large,column sep=large]
U_J \arrow[r] \arrow[d,"\unsure{\varepsilon|i}"'] & U \arrow[d] \\
X_J \arrow[r] & \operatorname{Spec}K=\xi \\
D_J \arrow[u,"\text{équiv. d'homotopie}"] & {}
\end{tikzcd}
}
d'où un hom. de suites exactes d'homotopie, \deleted{relative}\inserted{de celle relative} à la première \deleted{ligne}\inserted{colonne}, dans celle relative à la seconde
\diagram{
\begin{tikzcd}[row sep=huge,column sep=small]
\ldots \arrow[r,dash] &
\overset{\substack{\pi_2(D_J)\\\Vert}}{\pi_2(X_J)} \arrow[r] \arrow[d] &
\overset{\substack{T^J\\\Vert}}{\pi_1(\unsure{\bar G_m^{\times J}})} \arrow[r] \arrow[d] &
\overset{\substack{\mathcal E_{U_J}\\\Vert}}{\pi_1(U_J)} \arrow[r] \arrow[d] &
\overset{\substack{\mathcal E_{D_J}\\\Vert}}{\pi_1(X_J)} \arrow[r] \arrow[d] & 1 \\
\unsure{\ldots|--} \arrow[r,dash] &
\underset{\substack{\Vert\\0}}{\pi_2(\xi)} \arrow[r] &
\underset{\substack{\Vert\\\pi_{U/K}}}{\pi_1(U_{\bar\xi})} \arrow[r] &
\underset{\substack{\Vert\\\mathcal E_U}}{\pi_1(U)} \arrow[r,phantom,"\unsure{-|\to}" description] &
\underset{\substack{\Vert\\\mathcal E_{\unsure{K}}}}{\pi_1(\xi)} \arrow[r] & 1
\end{tikzcd}
}
qui montre que \inserted{le composé} $\pi_2(D_J)\longrightarrow T^J\longrightarrow\pi_{U/K}$ est nul --- donc si on admet que $T^J\unsure{-|\to}\pi_{U/K}$ est injectif, on conclura que $\pi_2(D_J)\unsure{-|\to}T^J$ est nul!

\margin{133}
